\chapter{Проектирование и реализация системы}

\section{Техническое задание}

\subsection{Назначение разработки}

Система BioSonification предназначена для генерации символической музыки на основе биологических последовательностей. Система должна преобразовывать FASTA-файлы в двухдорожечные MIDI-композиции, используя биологические признаки как управляющий сигнал.

\subsection{Функциональные требования}

\begin{enumerate}
\item Чтение FASTA-файлов с DNA, RNA и protein последовательностями
\item Фрагментация длинных последовательностей
\item Извлечение биологических признаков (состав, энтропия, белковые свойства)
\item Загрузка и структурирование MIDI-корпуса
\item Обучение модели гармонии
\item Обучение модели мелодии
\item Генерация двухдорожечного MIDI
\item Сохранение метаданных генерации
\end{enumerate}

\subsection{Нефункциональные требования}

\begin{itemize}
\item Обучение на GPU NVIDIA RTX 2060 6GB
\item Поддержка mixed precision для экономии памяти
\item Воспроизводимость результатов (фиксированный seed)
\item Сохранение checkpoint для повторного использования
\end{itemize}

\section{Архитектура системы}

Система построена по модульному принципу и включает следующие компоненты:

\begin{itemize}
\item \textbf{Биологический энкодер} — извлечение признаков из FASTA
\item \textbf{Музыкальный процессор} — структурирование MIDI-корпуса
\item \textbf{Pairing модуль} — сопоставление биологических и музыкальных сегментов
\item \textbf{Harmony model} — генерация аккордовой сетки
\item \textbf{Melody model} — генерация мелодии поверх гармонии
\item \textbf{MIDI renderer} — создание выходного файла
\end{itemize}

Общая схема пайплайна представлена на рисунке~\ref{fig:architecture}. Схема подчёркивает, что биологическая последовательность не преобразуется в ноты напрямую: сначала из неё извлекаются признаки, затем они используются для подбора обучающих пар и условной генерации.

\begin{figure}[htbp]
\centering
\includegraphics[width=0.95\textwidth]{pipeline_architecture.png}
\caption{Архитектура системы BioSonification}
\label{fig:architecture}
\end{figure}

Иерархия генерации показана на рисунке~\ref{fig:hierarchical-generation}. Биологический энкодер формирует два типа управляющей информации: компактный \texttt{control\_profile}, который попадает в префикс токенов, и 256-мерный embedding, который используется моделью как memory tokens.

\begin{figure}[htbp]
\centering
\includegraphics[width=0.95\textwidth]{hierarchical_generation_scheme.png}
\caption{Иерархическая генерация Bio→Harmony→Melody}
\label{fig:hierarchical-generation}
\end{figure}

\section{Данные и разбиение}

Для обучения использовались два типа данных: биологические последовательности в формате FASTA и музыкальные данные в формате MIDI. Биологическая часть формировалась из геномных данных RefSeq, музыкальная часть — из корпуса POP909. Корпус POP909 выбран как основной музыкальный источник, потому что он содержит структурированные MIDI-композиции с выраженной поп-гармонией и мелодическими партиями.

Длинные биологические последовательности разбивались на фрагменты фиксированной длины. Такой подход решает две задачи: ограничивает размер входа для энкодера и позволяет сопоставлять биологические участки с короткими музыкальными сегментами. Музыкальные композиции, в свою очередь, делились на 4-тактовые сегменты. Длина в 4 такта выбрана как компромисс: сегмент уже содержит локальную гармоническую структуру, но остаётся достаточно коротким для устойчивого обучения на GPU с 6GB видеопамяти.

После извлечения признаков и построения музыкальных сегментов выполнялось разбиение на train, validation и test выборки с фиксированным seed. Итоговая конфигурация актуальной модели содержит:
\begin{itemize}
\item 3746 биологических последовательностей и фрагментов;
\item 25940 музыкальных сегментов POP909;
\item 15925 обучающих пар;
\item 1870 validation-пар;
\item 935 test-пар.
\end{itemize}

Обзор данных и разбиения приведён на рисунке~\ref{fig:dataset-split}. Такое разбиение используется не только для обучения, но и для последующей интерпретации loss: validation и test loss считаются на парах, которые не использовались для обновления параметров модели.

\begin{figure}[htbp]
\centering
\includegraphics[width=0.95\textwidth]{dataset_and_split_overview.png}
\caption{Обзор подготовленных данных и разбиения выборки}
\label{fig:dataset-split}
\end{figure}

\section{Биологический энкодер}

Модуль \texttt{bio\_music\_pipeline/v2/bio.py} реализует извлечение признаков из биологических последовательностей.

\subsection{Поддерживаемые типы последовательностей}

Система автоматически определяет тип последовательности:
\begin{itemize}
\item DNA — содержит A, C, G, T
\item RNA — содержит A, C, G, U
\item Protein — содержит аминокислоты
\end{itemize}

\subsection{Фрагментация}

Длинные последовательности разбиваются на фрагменты с параметрами:
\begin{itemize}
\item \texttt{fragment\_length} = 1800 bp
\item \texttt{fragment\_stride} = 900 bp
\item \texttt{max\_fragments\_per\_record} = 24
\end{itemize}

\subsection{Извлекаемые признаки}

\textbf{Нуклеотидные признаки:}
\begin{itemize}
\item Частоты оснований (A, C, G, T/U)
\item Энтропия последовательности
\item GC-content
\item GC/AT skew
\item Кодоновая энтропия
\item k-mer статистика (k=1,2,3)
\end{itemize}

Энтропия последовательности вычисляется по распределению символов:
\begin{equation}
H(s) = - \sum_{a \in A} p(a)\log_2 p(a),
\end{equation}
где \(A\) — алфавит последовательности, \(p(a)\) — относительная частота символа \(a\). Для нуклеотидных последовательностей также используются доля GC и skew-показатели:
\begin{equation}
GC = \frac{n_G + n_C}{n_A + n_C + n_G + n_T},
\end{equation}
\begin{equation}
GC\text{-}skew = \frac{n_G - n_C}{n_G + n_C + \varepsilon}, \quad
AT\text{-}skew = \frac{n_A - n_T}{n_A + n_T + \varepsilon}.
\end{equation}
Малое значение \(\varepsilon\) предотвращает деление на ноль для коротких или вырожденных фрагментов.

\textbf{Белковые признаки:}
\begin{itemize}
\item Перевод через longest ORF
\item Ароматичность (Biopython ProtParam)
\item Индекс нестабильности
\item Изоэлектрическая точка
\item GRAVY (гидрофобность)
\item Молекулярная масса
\item Состав аминокислот
\end{itemize}

\subsection{Выходные данные}

Энкодер возвращает структуру \texttt{BioEncodingResult}:
\begin{itemize}
\item \texttt{vector} — 256-мерный embedding
\item \texttt{control\_profile} — 6-мерный профиль (tempo, density, harmony change, register, complexity, mode)
\item \texttt{tonic\_pc\_hint} — подсказка тонального центра
\end{itemize}

\section{Музыкальное представление}

Модуль \texttt{structured\_music.py} реализует структурированное представление музыки.

\subsection{Выбор формата представления}

На предварительном этапе рассматривались разные музыкальные форматы: MusicLang~\cite{musiclang}, MIDI, нотные форматы MusicXML/MEI/LilyPond и токенизированные MIDI-like представления. MusicLang интересен как семантический язык музыкального описания: он позволяет работать с аккордами, тональностью и относительными нотами на уровне, близком к музыкальному мышлению. Однако для финальной реализации он оказался менее удобен как базовый формат обучения, поскольку требует дополнительной конвертации и жёсткого гармонического контекста.

Нотные форматы MusicXML, MEI и LilyPond хорошо подходят для публикации партитур и визуализации результата, но содержат много информации, не влияющей напрямую на обучение генеративной модели: графические элементы, правила нотной гравировки, детальные обозначения исполнения. Для задачи генерации по биологическим данным такая избыточность усложняет парсинг и увеличивает объём промежуточного представления.

Поэтому в итоговой системе выбран комбинированный подход. MIDI используется как универсальный событийный формат хранения, воспроизведения и обмена~\cite{midi}. Поверх MIDI строится собственное структурированное представление, разделяющее музыку на гармонию и мелодию. Это сохраняет практическую совместимость MIDI и одновременно даёт модели более музыкально осмысленные токены: аккордовые события, позиции в такте, относительные pitch classes и длительности.

\subsection{Структуры данных}

\textbf{HarmonyBar} — аккорд в такте:
\begin{itemize}
\item \texttt{bar\_index} — номер такта
\item \texttt{root\_pc} — основной тон (0-11)
\item \texttt{quality} — качество аккорда (maj, min, dim, aug, dom7, maj7, min7, sus2, sus4, power)
\item \texttt{hold} — продолжение предыдущего аккорда
\end{itemize}

\textbf{MelodyEvent} — событие мелодии:
\begin{itemize}
\item \texttt{bar\_index} — номер такта
\item \texttt{position\_step} — позиция в такте
\item \texttt{duration\_steps} — длительность
\item \texttt{relative\_pc} — высота относительно аккорда
\item \texttt{octave} — октава
\end{itemize}

\subsection{Токенизация}

Система использует два токенизатора:
\begin{itemize}
\item \texttt{HarmonyTokenizer} — кодирует последовательность аккордов
\item \texttt{MelodyTokenizer} — кодирует мелодические события с учётом гармонии
\end{itemize}

\section{Pairing биологических и музыкальных сегментов}

Модуль \texttt{structured\_pairing.py} сопоставляет биологические фрагменты с музыкальными сегментами через пространство дескрипторов.

\subsection{Музыкальные дескрипторы}

Для каждого музыкального сегмента вычисляются:
\begin{itemize}
\item Темп (BPM)
\item Плотность нот на такт
\item Частота смены аккордов
\item Средний регистр
\item Сложность гармонии
\item Модальность (мажор/минор)
\end{itemize}

\subsection{Калибровка}

Биологический \texttt{control\_profile} калибруется под распределение музыкального корпуса, чтобы значения попадали в реалистичный диапазон.

Калибровка переводит биологический профиль \(b\) в шкалу музыкальных дескрипторов:
\begin{equation}
\hat{b} = \mathrm{clip}\left(
\frac{b - \mu_b}{\sigma_b + \varepsilon}\sigma_m + \mu_m,\;0,\;1
\right),
\end{equation}
где \(\mu_b, \sigma_b\) — среднее и стандартное отклонение биологических профилей, \(\mu_m, \sigma_m\) — параметры распределения музыкальных дескрипторов.

\subsection{Top-k matching}

Для каждого биологического фрагмента подбираются k=5 ближайших музыкальных сегментов по взвешенному евклидову расстоянию в пространстве дескрипторов:
\begin{equation}
d(\hat{b}, m) =
\sqrt{\sum_{i=1}^{6} w_i(\hat{b}_i - m_i)^2},
\end{equation}
где \(m\) — 6-мерный музыкальный дескриптор сегмента, \(w_i\) — веса отдельных признаков. Сегменты с меньшим расстоянием получают больший вес при обучении, потому что их музыкальные свойства лучше соответствуют биологическому профилю.

Логика сопоставления показана на рисунке~\ref{fig:descriptor-pairing}.

\begin{figure}[htbp]
\centering
\includegraphics[width=0.95\textwidth]{descriptor_pairing_scheme.png}
\caption{Сопоставление биологических и музыкальных сегментов в пространстве дескрипторов}
\label{fig:descriptor-pairing}
\end{figure}

\section{Нейросетевая модель}

Модуль \texttt{structured\_model.py} реализует \texttt{BioConditionedSequenceModel} — autoregressive Transformer с bio conditioning.

\subsection{Выбор архитектуры}

На этапе проектирования рассматривались два варианта: генеративно-состязательная архитектура с моделью-оценщиком и условная autoregressive-модель на основе Transformer~\cite{vaswani2017}. GAN-подход хорошо соответствует идее внутренней оценки результата, но для данной задачи оказался менее практичным: он требует устойчивого обучения двух моделей, сложнее воспроизводится на небольшом корпусе и хуже связывает численную функцию потерь с музыкальными свойствами результата.

Итоговая архитектура использует Transformer как обучаемую модель последовательностей. Её преимущество в том, что обучение сводится к предсказанию следующего токена по предыдущим токенам и биологическому условию. Это даёт интерпретируемую функцию потерь, позволяет отдельно обучать гармонию и мелодию, а качество результата проверять внешними метриками. Таким образом, роль ``оценщика'' перенесена из обучающего цикла в экспериментальную оценку: модель не содержит дискриминатор, но после генерации проверяется по valid MIDI rate, chord-tone ratio и другим дескрипторам.

\subsection{Архитектура}

Параметры актуальной модели:
\begin{itemize}
\item \texttt{d\_model} = 384
\item \texttt{n\_layers} = 6
\item \texttt{n\_heads} = 6
\item \texttt{dim\_feedforward} = 1536
\item \texttt{dropout} = 0.20
\item \texttt{bio\_memory\_tokens} = 4
\end{itemize}

\subsection{Bio conditioning}

Биологический вектор проецируется в пространство модели через MLP и представляется как memory tokens, которые участвуют в cross-attention.

Формально модель оценивает вероятность следующего токена с учётом предыдущей последовательности \(y_{<t}\) и биологического вектора \(z\):
\begin{equation}
P(y_t \mid y_{<t}, z) = \mathrm{softmax}(f_\theta(y_{<t}, z)).
\end{equation}
Для гармонической модели \(y\) является последовательностью аккордовых токенов, для мелодической модели — последовательностью мелодических событий с добавленным harmony prefix.

\subsection{Использование модели}

Одна и та же архитектура используется дважды:
\begin{enumerate}
\item Как harmony model — генерирует последовательность аккордов
\item Как melody model — генерирует мелодию с harmony prefix
\end{enumerate}

Обучение выполняется по функции cross-entropy:
\begin{equation}
\mathcal{L}(\theta) =
- \frac{1}{N}\sum_{n=1}^{N}\sum_{t=1}^{T_n}
\log P_\theta(y_{n,t} \mid y_{n,<t}, z_n).
\end{equation}
Префиксные токены, содержащие управляющий профиль и гармонический контекст, не используются как целевые токены loss; модель оптимизируется на предсказании музыкального продолжения.

\section{Процесс обучения}

\subsection{Конфигурация}

Актуальная модель использует конфиг \texttt{pipeline\_v2\_medium\_rtx2060\_fast.json}:
\begin{itemize}
\item Batch size: 8
\item Gradient accumulation: 4 (effective batch size 32)
\item Learning rate: 0.0003
\item Weight decay: 0.05
\item Mixed precision: enabled
\item Optimizer: AdamW
\item Segment length: 4 такта
\end{itemize}

\subsection{Параметры выборки}

\begin{itemize}
\item Bio sequences/fragments: 3746
\item Music segments: 25940 (POP909)
\item Train pairs: 15925
\item Validation pairs: 1870
\item Test pairs: 935
\end{itemize}

Эти значения соответствуют данным, описанным в разделе о подготовке выборки. В процессе обучения используются не исходные FASTA и MIDI напрямую, а подготовленные пары: биологический embedding, управляющий профиль, аккордовые токены и мелодические токены.

\subsection{Обучение harmony model}

\begin{itemize}
\item Epochs: 10
\item Best validation loss: 0.1454 (epoch 2)
\item Test loss: 0.1389
\item Training time: около 21 минуты
\end{itemize}

\subsection{Обучение melody model}

\begin{itemize}
\item Epochs: 15
\item Best validation loss: 0.1566 (epoch 7)
\item Test loss: 0.1650
\item Training time: около 31 минуты
\end{itemize}

\section{Реализация}

Проект реализован на Python 3.10+ с использованием:
\begin{itemize}
\item PyTorch 2.0+ — нейросетевой фреймворк~\cite{pytorch}
\item Biopython — биоинформатические инструменты~\cite{cock2009}
\item music21 — музыкальный анализ~\cite{music21}
\item mido — работа с MIDI
\item Flask — веб-интерфейс
\end{itemize}

Основные модули:
\begin{itemize}
\item \texttt{train\_bio\_music\_v2.py} — CLI обучения
\item \texttt{generate\_from\_fasta\_v2\_fragmented.py} — CLI генерации
\item \texttt{bio\_music\_pipeline/v2/} — основная логика
\item \texttt{web/app.py} — веб-интерфейс
\end{itemize}

Исходный код проекта размещён в открытом репозитории BioSonification/biosonification~\cite{biosonification_repo}. Более подробное описание программы, пользовательский сценарий запуска и сведения для разработчика приведены в приложении А.

\section*{Выводы по главе 2}

Во второй главе описано проектирование и реализация системы BioSonification. Разработана иерархическая архитектура Bio→Harmony→Melody. Обоснован переход от первоначально рассматриваемой генеративно-состязательной схемы к условному Transformer, поскольку выбранный подход устойчивее обучается на доступных данных и даёт более воспроизводимую процедуру оценки. Реализованы модули извлечения биологических признаков, структурирования музыки, pairing и нейросетевой генерации.

Система обучена на корпусе POP909 и геномных данных RefSeq. Актуальная модель показала validation loss 0.1454 для гармонии и 0.1566 для мелодии. Обучение выполнено на GPU RTX 2060 6GB с использованием mixed precision.
